% Generated by tools/edn_latex.py from reports/edn.md.
% Do not edit: change reports/edn.md and run `make edn`.
\ednentry{
  id = {EDN-13},
  title = {Keep the \texorpdfstring{\texttt{/\allowbreak{}feedback}}{/feedback} endpoint, but reachable only from inside the Compose network},
  date = {2026-09-23},
  milestone = {M6: Monitoring},
  topic = {API Design, ML System Design, Monitoring},
  participants = {@lukas2510},
  decision = {\texttt{POST /\allowbreak{}feedback} (FR-10) stays in the component and is no longer \texttt{[open]}. It takes a request ID and an observed price, stored as a delayed label, and FR-14 joins those labels to the prediction log to report the real error (MdAPE) and the empirical interval coverage per window, not only input drift. The endpoint is not publicly reachable: a reverse proxy in front of the API refuses \texttt{/\allowbreak{}feedback} from outside, so only the replay job inside the Compose network reaches it; the \texttt{X-\allowbreak{}API-\allowbreak{}Key} check stays behind that as a second layer in case the proxy rule is ever wrong. Because the key therefore never crosses the public internet, NFR-09 no longer requires TLS, and the previously \texttt{[proposed]} TLS item is dropped.},
  rationale = {\emph{Alternatives considered:}
\begin{itemize}[leftmargin=*, topsep=2pt]
\item \textbf{Option A: keep \texttt{/\allowbreak{}feedback} as a publicly reachable endpoint, protected by the API key alone.} The variant the requirements page described before this decision.
\newline Pros: one fewer component in the deployment; the endpoint behaves like a real public feedback channel.
\newline Cons: the key would travel in cleartext unless the VM gets TLS, and the FIB Virtech VM sits behind NAT with a forwarded port, so an HTTP-01 Let's Encrypt challenge on port 80 of a domain we control is most likely not possible; that left the whole point of adding the key resting on an unverified assumption about the VM.
\item \textbf{Option B: drop \texttt{/\allowbreak{}feedback}; the drift job reads the \texttt{ES} prices directly from the versioned artefact and joins them to the prediction log by row ID.}
\newline Pros: no write path, no key, no TLS question, no shared store needed for feedback; least code.
\newline Cons: the labels never enter the system through an interface, so the delayed-label mechanism exists only inside one batch job; the architecture no longer shows how ground truth would arrive in operation, which is the part M6 actually grades (``model performance degradation, and user impact''); EDN-03 chose the \texttt{ES} holdout precisely because every held-out listing has a price, and that argument only pays off if those prices flow back in somewhere.
\item \textbf{Option C (chosen): keep the endpoint, but make it unreachable from outside (reverse-proxy rule), with the API key as a second layer.}
\newline Pros: keeps the feedback loop as a real interface, so FR-14 can report error and coverage and FR-15's retraining is triggered by a measured degradation rather than by an input-distribution change alone; removes the TLS dependency entirely instead of leaving it pending on a VM capability nobody has confirmed; the attack surface of the only writing endpoint drops to zero from the internet; the reverse proxy is a component the course demo already suggests (nginx in front of \texttt{uvicorn}).
\newline Cons: one more component in \texttt{docker compose} and one proxy rule that has to be right, which is why the key is kept behind it and the deployment smoke test asserts from outside the VM that \texttt{/\allowbreak{}feedback} is refused.
\end{itemize}
\par
\emph{Why:} The endpoint's existence could not stay \texttt{[open]}, because FR-14, FR-15, NFR-03 and NFR-09 all depended on it. On whether to keep it, the decisive argument is what M6 is graded on: without labels the monitoring can only show that the inputs changed, while with them it can show that the error rises and the nominal 90\,\% intervals lose their coverage, which is also the stronger trigger for FR-15's retraining. The marginal cost is small, since the storage that dominates the work is needed for FR-13's prediction log anyway. On exposure, option C removes an unverified assumption (TLS on the Virtech VM) rather than carrying it as a pending decision, and it keeps the honest framing that traffic and labels are replayed from the \texttt{ES} holdout, not produced by real users.},
  ai = {Information seeking, Alternative generation, Alternative assessment, Recommendation},
  response = {Accepted},
  assessment = {During the review of PR \#19, AI flagged that FR-10 carried \texttt{[open]} and \texttt{[decided]} markers in the same row and that four other requirements silently depended on it, then laid out options A to C with the M6 rubric wording and the likely NAT setup of the Virtech VM as the deciding factors, and recommended keeping the endpoint while making it internal. Asked to explain the recommendation in more detail, AI also named the counter-argument itself (the labels are simulated, so a critic could call the endpoint a facade) and the mitigation (state plainly in the report that traffic and labels are replayed). Lukas reviewed the reasoning and accepted the recommendation.},
  aievidence = {Claude Code session on 2026-09-23 while reviewing PR \#19: AI recommended keeping \texttt{/\allowbreak{}feedback} and combining it with internal-only exposure; Lukas asked ``erkläre mir FR-10 genauer und warum du das eine vorschlägst'', AI explained the delayed-label mechanism, the dependencies and the three options; Lukas replied ``ok ich mag deine empfehlung ändere es so''.},
  evidence = {\href{https://github.com/mlops-2627q1-mds-upc/MLOPS_Recommenditos/blob/main/docs/docs/specification.md}{Specification} FR-10, FR-14, NFR-09; \href{https://github.com/mlops-2627q1-mds-upc/MLOPS_Recommenditos/blob/main/docs/docs/project-brief.md}{project brief} sections 3.2 and 5; \hyperref[EDN-03]{EDN-03}, \hyperref[EDN-12]{EDN-12}; \href{https://github.com/mlops-2627q1-mds-upc/MLOPS_Recommenditos/pull/19}{PR \#19}.},
}
